\documentclass[12pt, a4paper]{article}

\usepackage[utf8]{inputenc}
\usepackage[italian]{babel}
\usepackage[T1]{fontenc}
\usepackage{lmodern}
\usepackage{microtype}
\usepackage{amsmath, amssymb}
\usepackage{geometry}
\usepackage{enumitem}
\usepackage{booktabs}
\usepackage{eurosym}
\usepackage{graphicx}
\usepackage{subcaption}
\usepackage{titlesec}
\usepackage{tcolorbox}
\usepackage{hyphenat}
\usepackage{tocloft}
\usepackage{xcolor}
\usepackage{hyperref}
\usepackage{datetime}
\usepackage{array}
\usepackage{float}
\usepackage{listings}
\usepackage{longtable}
\usepackage{pdflscape}
\usepackage{fancyhdr}
\usepackage{refcount}

% Colonna di tabella allineata a sinistra, senza giustificazione forzata.
\newcolumntype{L}[1]{>{\raggedright\arraybackslash}p{#1}}

% --- CONFIGURAZIONE LISTINGS ---
\definecolor{codegray}{rgb}{0.5,0.5,0.5}
\definecolor{codegreen}{rgb}{0.1,0.5,0.1}
\definecolor{codeblue}{rgb}{0.1,0.1,0.6}
\lstset{
  basicstyle=\ttfamily\small,
  keywordstyle=\color{codeblue}\bfseries,
  commentstyle=\color{codegreen}\itshape,
  stringstyle=\color{codegray},
  frame=single,
  breaklines=true,
  showstringspaces=false,
  captionpos=b,
  columns=fullflexible,
  keepspaces=true
}

% --- CONTATORI E COMANDI PER REQUISITI, USER STORY E USER FLOW ---
\makeatletter

% Livelli personalizzati per l'indice
\def\toclevel@RF{2}
\def\toclevel@subRF{3}
\def\toclevel@RNF{2}
\def\toclevel@subRNF{3}
\def\toclevel@US{2}
\def\toclevel@UF{2}

\newcommand{\l@RF}[2]{\begingroup\small\@dottedtocline{2}{1.5em}{2.3em}{#1}{#2}\endgroup}
\newcommand{\l@subRF}[2]{\begingroup\small\@dottedtocline{3}{3.8em}{3.4em}{#1}{#2}\endgroup}
\newcommand{\l@RNF}[2]{\begingroup\small\@dottedtocline{2}{1.5em}{2.3em}{#1}{#2}\endgroup}
\newcommand{\l@subRNF}[2]{\begingroup\small\@dottedtocline{3}{3.8em}{3.4em}{#1}{#2}\endgroup}
\newcommand{\l@US}[2]{\begingroup\small\@dottedtocline{2}{1.5em}{2.3em}{#1}{#2}\endgroup}
\newcommand{\l@UF}[2]{\begingroup\small\@dottedtocline{2}{1.5em}{2.3em}{#1}{#2}\endgroup}

% --- Formato delle sezioni ---
\titleformat{\section}
  {\normalfont\LARGE\bfseries}
  {\thesection}
  {6pt}
  {}

% --- Requisiti funzionali (RF) ---
\newcounter{reqfunc}
\newcounter{subRFcounter}[reqfunc]
\renewcommand{\thesubRFcounter}{\arabic{reqfunc}.\arabic{subRFcounter}}

% Numero a due cifre (RF-01, RNF-01, US-01, UF-01)
\newcommand{\twodigits}[1]{\ifnum#1<10 0\fi\number#1}

\newcommand{\RF}[1]{%
  \refstepcounter{reqfunc}%
  \setcounter{subRFcounter}{0}%
  \addcontentsline{toc}{RF}{RF-\twodigits{\arabic{reqfunc}}: #1}%
  \paragraph*{RF-\twodigits{\arabic{reqfunc}}: #1}%
}

\newcommand{\subRF}[1]{%
  \refstepcounter{subRFcounter}%
  \addcontentsline{toc}{subRF}{RF \thesubRFcounter: #1}%
  \noindent\hspace{1.5em}\textbf{RF \thesubRFcounter:} #1\par
}

% --- Requisiti non funzionali (RNF) ---
\newcounter{reqnonfunc}
\newcounter{subRNFcounter}[reqnonfunc]
\renewcommand{\thesubRNFcounter}{\arabic{reqnonfunc}.\arabic{subRNFcounter}}

\newcommand{\RNF}[1]{%
  \refstepcounter{reqnonfunc}%
  \setcounter{subRNFcounter}{0}%
  \addcontentsline{toc}{RNF}{RNF-\twodigits{\arabic{reqnonfunc}}: #1}%
  \paragraph*{RNF-\twodigits{\arabic{reqnonfunc}}: #1}%
}

\newcommand{\subRNF}[1]{%
  \refstepcounter{subRNFcounter}%
  \addcontentsline{toc}{subRNF}{RNF \thesubRNFcounter: #1}%
  \noindent\hspace{1.5em}\textbf{RNF \thesubRNFcounter:} #1\par
}

% --- User story (US) ---
\newcounter{userStories}

\newcommand{\US}[1]{%
  \refstepcounter{userStories}%
  \addcontentsline{toc}{US}{US-\twodigits{\arabic{userStories}}: #1}%
  \subsection*{US-\twodigits{\arabic{userStories}} --- #1}%
}

% --- User flow (UF) ---
\newcounter{userFlows}

\newcommand{\UF}[1]{%
  \refstepcounter{userFlows}%
  \addcontentsline{toc}{UF}{UF-\twodigits{\arabic{userFlows}}: #1}%
  \subsection*{UF-\twodigits{\arabic{userFlows}} --- #1}%
}

\makeatother

% Riferimenti nel formato stabile RF-01, RNF-01, US-01, UF-01.
\newcommand{\RFref}[1]{RF-\twodigits{\getrefnumber{#1}}}
\newcommand{\RNFref}[1]{RNF-\twodigits{\getrefnumber{#1}}}
\newcommand{\USref}[1]{US-\twodigits{\getrefnumber{#1}}}
\newcommand{\UFref}[1]{UF-\twodigits{\getrefnumber{#1}}}

\newcommand{\myfootertext}{}
\newcommand{\myicon}{logo_unitn_chiaro.png}
\newcommand{\myappicon}{logo_app.png}

\definecolor{lightergray}{gray}{0.35}
\definecolor{tocsec}{gray}{0.35}
\definecolor{tocsubsec}{gray}{0.45}
\pagecolor{white}
\color{black}

\hypersetup{
    colorlinks=true,
    linkcolor=lightergray,
    urlcolor=cyan,
    pdftitle={TEPO --- D2},
    pdfauthor={Gruppo T14},
    pdfsubject={Documento di Requirements Engineering --- baseline},
    pdfkeywords={TEPO, requisiti, use case, user story, energy portfolio optimization}
}

\renewcommand{\cftsecfont}{\color{tocsec}\bfseries}
\renewcommand{\cftsecpagefont}{\color{tocsec}\bfseries}
\renewcommand{\cftsubsecfont}{\color{tocsubsec}}
\renewcommand{\cftsubsecpagefont}{\color{tocsubsec}}
\renewcommand{\cftsubsecleader}{\color{tocsubsec}\cftdotfill{\cftdotsep}}

\geometry{margin=2.5cm,headheight=24.5pt}
\setlength{\parskip}{0.35em}
\renewcommand{\arraystretch}{1.2}

% --- Intestazione e piè di pagina ---
\pagestyle{fancy}
\fancyhf{}
\fancyhead[L]{\small\color{tocsubsec}TEPO --- D1}
\fancyhead[R]{\small\color{tocsubsec}Gruppo T14}
\fancyfoot[C]{\small\thepage}
\renewcommand{\headrulewidth}{0.4pt}
\renewcommand{\headrule}{\hbox to\headwidth{\color{tocsubsec}\leaders\hrule height \headrulewidth\hfill}}
\fancypagestyle{plain}{\fancyhf{}\fancyfoot[C]{\small\thepage}\renewcommand{\headrulewidth}{0pt}}

\begin{document}

\begin{titlepage}
    \thispagestyle{empty}
    % --- Loghi in alto (centrati, stessa riga) ---
    \begin{center}
        \href{https://www.unitn.it/}{%
            \includegraphics[height=3.5cm, keepaspectratio]{\myicon}%
        }%
        \hspace{1.5cm}%
        \href{https://github.com/T-E-P-O/D2}{%
            \includegraphics[height=3.5cm, keepaspectratio]{\myappicon}%
        }%
    \end{center}

    \vspace{1.2cm}

    % --- Titolo e Dipartimento (Allineati a sinistra) ---
    \begin{flushleft}
        {\scshape\Large Università degli Studi di Trento \par}
        {\large Dipartimento di Ingegneria e Scienza dell'Informazione \par}

        \vspace{1.3cm}

        {\Huge\bfseries TEPO --- D2 \par}
        \vspace{0.3cm}
        {\Large\itshape Trento Energy Portfolio Optimizer \par}
        \vspace{0.35cm}
        {\large Documento di Requirements Engineering --- baseline \par}
        \vspace{0.5cm}
        {\color{cyan}\rule{\textwidth}{2pt}}
    \end{flushleft}

    \vspace{0.7cm}

    % --- Document Info (come da template D1) ---
    \begin{flushleft}
        {\small
        \renewcommand{\arraystretch}{1.25}
        \begin{tabular}{@{}p{3.1cm}p{11.7cm}@{}}
            \textbf{Doc. Name} & TEPO --- D2, Descrizione di Progetto \\
            \textbf{Description} & Documento di Requirements Engineering: progetto, requisiti funzionali e non funzionali, use case, user story, acceptance criteria, mockup e traceability. \\
        \end{tabular}}
    \end{flushleft}

    \vspace{0.9cm}

    % --- Gruppo e autori ---
    \begin{flushleft}
        {\Large \textbf{Gruppo T14} \par}
        \vspace{0.3cm}
        \renewcommand{\arraystretch}{1.3}
        \begin{tabular}{@{}l}
            \Large \href{https://github.com/mirconegri}{Mirco Negri} \\
            \Large \href{https://github.com/ghebbb}{Gabriele Duca} \\
            \Large \href{https://github.com/scheddiee}{Riccardo Scheda} \\
        \end{tabular}
    \end{flushleft}

    \vfill

    % --- Piè di pagina ---
    \begin{flushleft}
        {\small Corso di Ingegneria del Software --- A.A. 2026/2027 \par}
        {\small Progetto nell'ambito dell'iniziativa ``100 progetti per il Comune di Trento'' \par}
        \vspace{0.2cm}
        {\small\textit{\myfootertext} \par}
        {\large \today \par}
    \end{flushleft}
\end{titlepage}

\newpage
\begingroup
\hypersetup{hidelinks}
\tableofcontents
\endgroup

\newpage

\include{sections/Il_nostro_progetto}
\include{sections/RF}
\include{sections/RNF}
\include{sections/UD}
\include{sections/US}
\include{sections/DesignFrontEnd}
\include{sections/UF}
\include{sections/Traceability}
\include{sections/GestioneModifiche}

\end{document}
